\chapter{Introduction}\label{ch:introduction}

\section{Motivation}

Biomimetic robots provide a platform for investigating how mechanical structure and control interact to generate movement, offering complementary insights to simulation-based studies \citep{andrews_general_1985, millard_flexing_2013, arnold_model_2010}. Wheeled robots remain the platform of choice in structured environments, where their speed, efficiency, and payload capacity are difficult to match; in the unstructured terrain of the natural and built world, legged machines hold the advantage, and biomimetic legs that reproduce human force capabilities have the potential to excel where wheels cannot \citep{frey_legged_2026, ha_learning_2025}. Musculoskeletal humanoids and limb systems driven by tendon-like actuators have been shown to support biologically meaningful experiments on motion generation, internal force transmission, and neural control hypotheses \citep{asano_design_2017, asano_musculoskeletal_2019, hitzmann_anthropomorphic_2018, liu_robotic_2018, chen_realizing_2019}. The dual perspective of using robots to learn about biology, and using biology to design more adaptive robots, has been emphasized in both biomimetic robotics and biomechanics \citep{ijspeert_amphibious_2020, asano_design_2017}. With such robotic systems, experiments can be performed that would not be practical or ethical with human test subjects, such as removing a specific muscle from the system or blocking a sensory feedback pathway. Experiments of this kind allow direct testing of how the components of the neuromechanical system interact, and they offer new opportunities for testing neural control theories that are faster and cheaper than human observation studies \citep{shin_constructive_2018, goldsmith_investigating_2021, schilling_adaptive_2022}.

Artificial muscles are a key enabling technology for this class of robot. Among the available actuation technologies, McKibben-style braided pneumatic actuators (BPAs) combine low mass, high force-to-weight ratio, and intrinsic compliance with a force--length curve that qualitatively resembles that of skeletal muscle \citep{liang_comparative_2020, hunt_modeling_2017}. BPAs have been used in quadrupedal robots \citep{aschenbeck_design_2006, hunt_modeling_2017, scharzenberger_design_2019}, musculoskeletal humanoids \citep{asano_design_2017, hitzmann_anthropomorphic_2018}, and dexterous manipulation systems \citep{chen_realizing_2019}. However, BPAs also differ from biological muscle in important ways: their maximum tensile force occurs at resting length rather than at an optimal fiber length, their maximum contraction ratio is substantially lower than that of human muscle, and their force output is a highly nonlinear function of internal pressure, contraction, and loading state \citep{hunt_modeling_2017, sarosi_comparative_2017, liang_comparative_2020}. The author's experience with \qty{10}{\mm} actuators made this gap concrete: manufacturer specifications overstate the force and contraction available at the short resting lengths a leg design requires, published models fitted at one cut length do not transfer to other resting lengths, and designs carried out against those specifications delivered less range of motion and joint torque than intended, forcing iteration on hardware rather than at the design stage. Designing a biomimetic robot around BPAs therefore requires quantitative, predictive models of the actuator and of the joint torques it can produce.

\section{The Biomimetic Humanoid Robot Project}\label{sec:project}

This dissertation is part of a long-running effort at Portland State University to build a biomimetic humanoid robot from the trunk down, actuated by pneumatic artificial muscles and controlled by a synthetic neural network (SNN). The project's guiding hypothesis is that a robot with human-like biomechanics --- human-like muscle numbers, architecture, and routing, human-like joint ranges of motion, and human-like isometric torque profiles --- controlled by a biologically grounded neural architecture, will be capable of biomimetic movement and will serve as a testbed for neuromechanical hypotheses, including loss-of-function experiments that are impossible in living animals. A by-product of this work is expected to be improved prosthetic and orthotic device design.

The present work builds directly on a series of studies by the author and collaborators. In the author's master's project, the kinematics and dynamics of a biomimetic robot leg were matched against the Gait2392 OpenSim model, artificial muscle attachment and wrapping locations were determined for the robot legs, and MATLAB tooling was developed for evaluating muscle lengths, moment arms, and isometric joint torques over the full range of motion \citep{bolen_bipedal_2019, bolen_determination_2019}. That work showed that simply transplanting human muscle attachment locations onto a BPA-actuated robot does not reproduce human torque profiles; the placement optimization that closes this gap was completed during this dissertation and is summarized with the results in Section~\ref{sec:foundation}. In parallel, a biomimetic four-bar knee linkage with a migrating instantaneous center of rotation was designed and validated by Steele \citep{steele_development_2017, steele_biomimetic_2018, steele_biomimetic_2018-1}, and a canine-inspired quadruped platform (Muscle Mutt, formerly  called DoggyDeux) with antagonistic BPA pairs and synthetic neural network control was developed by \citet{scharzenberger_design_2019}, extending an earlier demonstration by the Agile and Adaptive Robotics Laboratory (AARL) of central pattern generator (CPG) control of a pneumatically actuated robot \citep{hunt_development_2017}.

\section{Problem Statement}\label{sec:problem}

Within this context, three gaps stand between the existing robot design and a robot that can walk under neural control.

First, the actuator force models available at the start of this work were inadequate for design purposes. The manufacturer's simulation tool over-predicts the force of short BPAs by large margins, and the published empirical models are fitted to a single resting length each, so they do not accurately predict how force changes when a designer chooses a different actuator length \citep{sarosi_comparative_2017, martens_modeling_2017, hunt_modeling_2017}. Because muscle resting length is one of the primary degrees of freedom available to the robot designer, this gap blocks rational actuator selection.

Second, actuator force alone does not determine joint torque. The robot knee applies BPA force through brackets, artificial tendons, and wrapping paths. It was unknown how much these effects reduce usable torque relative to the rigid-body prediction, or how to best correct for them. Additionally, the biomimetic knee's optimized muscle routing \citep{morrow_optimization_2020} had never been validated experimentally or compared against human isometric torque data.

Third, no simulation pipeline existed for closing the loop between a musculoskeletal robot model of this kind and a biologically realistic neural controller. The Neuromechanical Modeling literature provides the ingredients --- two-level spinal CPG architectures \citep{rybak_modelling_2006, mccrea_organization_2008}, neuromechanical simulation frameworks \citep{markin_afferent_2010}, and synthetic nervous system tooling \citep{nourse_sns_2023} --- but they had not been assembled for a BPA-driven bipedal robot with human-like routing.

\section{Research Objectives}\label{sec:objectives}

This dissertation addresses the first two gaps experimentally and lays the groundwork for the third. The objectives are:

\begin{enumerate}
    \item \textbf{Characterize the isometric force of Festo BPAs} with uninflated diameters of \qtylist{10;20}{\mm} over a wide range of resting lengths, contractions, and pressures; derive improved maximum-force and normalized force--strain--pressure models; and compare them against existing models (Chapters~\ref{ch:methods} and~\ref{ch:results}, Section~\ref{sec:results_force}).
    \item \textbf{Validate and improve the prediction of isometric knee torque} produced by BPAs on a pinned joint and on a biomimetic four-bar knee; identify and model the losses due to constant length offsets, series compliance, and wrapping; and compare the resulting torque envelopes with human capabilities (Section~\ref{sec:results_torque}).
    \item \textbf{Define a path to neuromechanical control} of the resulting robot using a converted MuJoCo model, the SNS-Toolbox synthetic nervous system environment, and a two-layer CPG architecture with proprioceptive and contact feedback (Chapter~\ref{ch:futurework}).
\end{enumerate}

\section{Dissertation Organization}\label{sec:organization}

Chapter~\ref{ch:background} reviews the relevant background: biomimetic and musculoskeletal robots, artificial muscle technologies and BPA modeling, human lower-limb biomechanics and the OpenSim benchmark models, the neural control of locomotion with an emphasis on two-level CPG architectures and sensory afferent feedback, synthetic nervous systems and the Sensory Afferent Database, and the simulation tooling that connects OpenSim models to MuJoCo. Chapter~\ref{ch:methods} presents the materials and methods: the actuator force model and its characterization experiment, the pinned-knee and biomimetic knee test stands, the hybrid torque calculation method that isolates the loss mechanisms, the correction terms fitted to them, and the optimization algorithms used to redesign actuator placements. Chapter~\ref{ch:results} first summarizes the completed foundation studies (the placement optimization, the Sensory Afferent Database, the BPA balance platform, and the BPA pressure-dynamics model), then reports the actuator characterization, the identification of the torque correction terms on the pinned and biomimetic knees, the redesigned actuator routes, and the experimental comparison of the biomimetic knee with human isometric torque. Chapter~\ref{ch:discussion} discusses the findings, their limitations, and their implications for real-time control. Chapter~\ref{ch:futurework} assembles the validated plant and the neural controller into a working neuromechanical simulation pipeline, an OpenSim-to-MuJoCo conversion coupled to a two-layer CPG with full sensory feedback; with these tools in place, researchers can run lesion, stimulation, and gait-modulation experiments on the simulated humanoid leg today, and the same tools define the path to a treadmill robot with embedded real-time control. Chapter~\ref{ch:conclusion} concludes.
